\section{Setup, notation, and three levels}
\label{p:setup}

\subsection{Notation}
We use the conventions established in the long form of this work:
$\vk$ denotes a Bloch momentum on the discrete model, $\vq$ a small
deviation about a band-crossing point $\vk_{0}$, and $\vp$ the continuum
infrared momentum. Internal degrees of freedom are organized by
$\sigma_{i}$ (pseudospin / Weyl two-component) and $\tau_{i}$ (chiral
index or coupled-block label). The canonical four-component Dirac block
that we recover at the end of the ladder takes the form
\begin{equation}
\HD \;=\; v\,\tau_{z} \otimes (\sigma_{x} p_{x} + \sigma_{y} p_{y} + \sigma_{z} p_{z})
            \;+\; m\,\tau_{x}\otimes\bbI ,
\label{p:eq:HDcanon}
\end{equation}
whose squared form
$\HD^{2} = (v^{2}\lvert \vp\rvert^{2} + m^{2})\,\bbI_{4}$
will reappear several times.

\subsection{Three levels}
A recurring source of confusion in this corner of the literature is the
conflation of three distinct levels of statement. We commit to keeping
them apart throughout the paper:

\begin{description}
\item[\textbf{Level I (microscopic exact).}]
    The discrete substrate: lattice or step protocol $\Uop$, locality, the
    exact group of discrete symmetries, quasi-energy periodicity in
    discrete-time models. At this level the fundamental dynamical object
    is $\Uop$, not $\Hop$.
\item[\textbf{Level II (effective infrared).}]
    The linearization around a band-crossing point, the appearance of
    Weyl or Dirac blocks, approximate continuous isotropy, the reading
    of mass, chirality, and effective velocity. Here it is natural to
    write $\Uop \simeq \bbI - \ii\,\Delta t\,\Hop_{\mathrm{eff}}$.
\item[\textbf{Level III (covariant spinorial).}]
    The Clifford algebra of $\gmu$, the spinorial Lorentz generators
    $\sigmunu$, the proper orthochronous group $SO^{+}(3,1)$,
    Poincar\'e, and the double-cover $SL(2,\bbC) \to SO^{+}(3,1)$. This
    level is the \emph{organization} of the infrared limit, not an exact
    symmetry of the substrate.
\end{description}

We will reserve the adjective \enquote{exact} for Level I, and the
adjective \enquote{effective} or \enquote{infrared} for Levels II and
III. This split, banal as it sounds, is what saves the rest of the
construction from over-claiming.

\subsection{The protospace tuple}
By \emph{protospace} we mean the ampliated structure that carries the
microscopic dynamics:
\begin{equation}
\Proto \;\sim\; \bigl(\Lattice,\,\bbC^{N},\,\Uop,\,G_{\mathrm{micro}}\bigr) ,
\label{p:eq:proto-tuple}
\end{equation}
where $\Lattice$ is the underlying graph (not necessarily periodic),
$\bbC^{N}$ the local internal Hilbert space, $\Uop$ the local step, and
$G_{\mathrm{micro}}$ its exact discrete symmetry group. Identifying
$\Proto$ with the Brillouin torus $T^{d}$ would confuse a
\emph{calculational device} (the Fourier-space representation of a
periodic protocol) with the \emph{underlying structure}. We come back to
this point in Section~\ref{p:open-problems}.
